\documentclass[a4paper,11pt]{article}
\usepackage{amsmath,amssymb}
\usepackage{booktabs}
\usepackage{lmodern}
\usepackage[utf8]{inputenc}
\usepackage[T1]{fontenc}
\usepackage[french]{babel}
\usepackage{float}
\usepackage{graphicx}
\usepackage{geometry}
\usepackage{enumitem}
\usepackage[hidelinks]{hyperref}
\geometry{margin=2.5cm}
\newcommand{\fig}[2][0.9\textwidth]{\IfFileExists{Figures/#2}{\includegraphics[width=#1]{Figures/#2}}%
  {\fbox{\parbox[c][3cm][c]{#1}{\centering\small\texttt{\detokenize{Figures/#2}}}}}}

\begin{document}
\begin{titlepage}
    \begin{center}
        \vspace*{1cm}
            
        \Huge
        \textbf{Optimisation}
            
        \vspace{0.5cm}
        \LARGE
        Projet 6 -- Algorithme EM \\
        Estimation d'un mélange de Gaussiennes
        
        \vspace{0.5cm}

        \large Professeur : Stéphane CHRÉTIEN
            
        \vspace{1.5cm}
            
        \textbf{Eoin BRERETON HURLEY}

        \vspace{0.2cm}
        
        \large \today

        \vspace{0.2cm}
            
        \vspace{2.0cm}
            
        \includegraphics[width=1.0\textwidth]{Figures/logo ICOM.png}

        \includegraphics[width=0.8\textwidth]{Figures/logo LYON 2.png}
        
        \vspace{1.0cm}
            
    \end{center}
\end{titlepage}

\section{Introduction}
Nous reproduisons les expériences de la section EM de la leçon~19 d'EE~227C~\cite{ee227c} : l'estimation d'un mélange
de deux Gaussiennes en dimension~1, puis la comparaison avec \emph{k-means}. Nous étudions ensuite l'influence de
l'initialisation, de la séparation des composantes, de la taille d'échantillon et du nombre de composantes, avant une
application aux données Iris.

\section{L'algorithme EM}
Soit $z_i\in\{1,\dots,K\}$ une variable latente avec $P(z_i=k)=\pi_k$ et $x_i\mid z_i=k\sim\mathcal N(\mu_k,\sigma_k^2)$.
On cherche le maximum de vraisemblance de $\theta=(\pi,\mu,\sigma)$ :
\[
L(\theta)=\sum_{i=1}^n\log\sum_{k=1}^K\pi_k\,\varphi(x_i;\mu_k,\sigma_k^2).
\]
$L$ n'est pas concave et n'a pas de maximiseur explicite. EM~\cite{ssbd} alterne deux étapes explicites :
\begin{itemize}[nosep]
\item \textbf{E :} $t_{ik}=\dfrac{\pi_k\varphi(x_i;\mu_k,\sigma_k^2)}{\sum_j\pi_j\varphi(x_i;\mu_j,\sigma_j^2)}$,
  calculé en échelle logarithmique (\emph{log-sum-exp}) ;
\item \textbf{M :} avec $n_k=\sum_it_{ik}$, on pose $\pi_k=n_k/n$, $\mu_k=\sum_it_{ik}x_i/n_k$ et
  $\sigma_k^2=\sum_it_{ik}(x_i-\mu_k)^2/n_k$.
\end{itemize}
EM est une maximisation alternée de
$G(Q,\theta)=\sum_{i,k}t_{ik}\log p_\theta(x_i,k)-\sum_{i,k}t_{ik}\log t_{ik}$, qui coïncide avec $L(\theta)$ après
l'étape~E. La log-vraisemblance ne décroît donc jamais, mais EM peut s'arrêter en un maximum local ou en un point selle.
Les dérivations ci-dessous (figures~\ref{fig:resp-em} et~\ref{fig:estim-em}) traitent le cas multivarié ; en
dimension~1, avec $\Sigma_k=\sigma_k^2$, on retrouve les formules ci-dessus.

\begin{figure}[H]\centering \fig[0.75\textwidth]{EM-Responsibility.jpg}
\caption{Calcul des Responsabilités, $\tau_{ik}$ ou $t_{ik}$}\label{fig:resp-em}\end{figure}

\begin{figure}[H]\centering \fig[0.75\textwidth]{EM-Estimator-Proofs-Multivariate.jpg}
\caption{Preuves des Estimateurs GMM (cas multivarié)}\label{fig:estim-em}\end{figure}

%\paragraph{Correction du code de Berkeley.} La fonction \texttt{gmm\_likelihood} de Berkeley renvoie
%\texttt{logsumexp(joint)} sans préciser l'axe. Elle calcule donc $\log\sum_i\sum_kp_\theta(x_i,k)$, et non
%$\sum_i\log\sum_kp_\theta(x_i,k)$. Avec cette version, le critère BIC de la section~4.4 choisit toujours $K=1$. Nous
%utilisons \texttt{logsumexp(joint, axis=1).sum()}.

\section{Expérience de Berkeley}
Les données sont $n=4000$ points, pour moitié $\mathcal N(-2,1)$ et pour moitié $\mathcal N(2,1)$ ; les moyennes
sont initialisées aux quartiles. En 20 itérations, EM estime $\pi=(0{,}498;0{,}502)$, $\mu=(-2{,}02;1{,}99)$ et
$\sigma=(0{,}997;1{,}002)$. La log-vraisemblance est croissante et se stabilise après environ 5 itérations
(figure~\ref{fig:berk}).
\begin{figure}[ht]\centering \fig{berkeley_em.png}
\caption{Log-vraisemblance au fil des itérations et densité estimée.}\label{fig:berk}\end{figure}

\section{Influence des paramètres}
\subsection{Initialisation}
Les 20 initialisations aléatoires (moyennes tirées parmi les données) atteignent toutes le même maximum
($L=-8205{,}58$), en 33 à 56 itérations, contre 35 pour les quartiles. Une initialisation symétrique
($\mu_1=\mu_2$, $\sigma_1=\sigma_2$, $\pi_1=\pi_2$) donne en revanche $t_{i1}=t_{i2}=\frac12$ pour tout $i$ : les deux
composantes restent identiques à chaque itération. EM s'arrête alors en un point selle ($L=-8902{,}4$, valeur
d'une seule Gaussienne) (figure~\ref{fig:init}).

\subsection{Séparation des composantes}
Avec des moyennes $\pm\Delta/2$, le nombre d'itérations nécessaires pour que $L$ varie de moins de $10^{-10}$ vaut 34
($\Delta=4$), 142 ($\Delta=3$), 1472 ($\Delta=2$) et plus de 2000 ($\Delta=1$). La convergence d'EM est linéaire, et
son taux dépend de la part d'information manquante : plus les composantes se recouvrent, plus les $t_{ik}$ sont
incertains et plus EM est lent (figure~\ref{fig:init}). Les moyennes restent bien estimées
($-0{,}58$ et $0{,}52$ pour $\Delta=1$), mais leur précision se dégrade quand $\Delta$ diminue.
\begin{figure}[ht]\centering \fig[0.48\textwidth]{initialisation.png}\hfill\fig[0.48\textwidth]{separation.png}
\caption{Gauche : 20 initialisations aléatoires (gris), quartiles et initialisation symétrique. Droite : écart à la
log-vraisemblance finale selon la séparation $\Delta$.}\label{fig:init}\end{figure}

\subsection{Taille d'échantillon et nombre de composantes}
L'erreur moyenne sur les moyennes, calculée sur 20 répétitions, décroît comme $1/\sqrt n$ entre $n=50$ et $n=10^4$
(figure~\ref{fig:n}). C'est l'erreur statistique de l'estimateur du maximum de vraisemblance, qu'aucune
optimisation ne peut réduire. La log-vraisemblance ne progresse presque plus au-delà de $K=2$ (de $-8205{,}6$ à
$-8204{,}6$), et le critère $\mathrm{BIC}=-2L+(3K-1)\log n$ retient bien $K=2$ :
\begin{center}\small
\begin{tabular}{lccccc}\toprule
$K$ & 1 & 2 & 3 & 4 & 5\\\midrule
BIC (simulé) & 17821 & \textbf{16453} & 16476 & 16500 & 16525\\
BIC (Iris) & 605 & \textbf{426} & 440 & 455 & 470\\\bottomrule
\end{tabular}
\end{center}
\begin{figure}[ht]\centering \fig[0.45\textwidth]{taille_echantillon.png}\hfill\fig[0.5\textwidth]{iris.png}
\caption{Gauche : erreur d'estimation en fonction de $n$. Droite : mélange retenu pour les pétales d'Iris.}\label{fig:n}\end{figure}

\section{\emph{k-means} et données réelles}
Sur l'exemple de J.~VanderPlas~\cite{vdp} (4 groupes en dimension~2), \emph{k-means} et le mélange de Gaussiennes ajusté par
EM donnent la même partition (figure~\ref{fig:km}). \emph{k-means} est la limite d'EM pour des Gaussiennes
sphériques de même variance $\sigma^2\to0$ : les $t_{ik}$ deviennent des affectations 0/1.

\paragraph{Iris.} Pour la longueur des pétales ($n=150$), le BIC retient $K=2$, avec $\pi=(0{,}33;0{,}67)$ et
$\mu=(1{,}46;4{,}91)$ cm, comme scikit-learn. Le premier groupe contient exactement les 50 \emph{setosa} ; le second
réunit \emph{versicolor} et \emph{virginica}, que cette seule variable ne sépare pas.
\begin{figure}[ht]\centering \fig[0.75\textwidth]{kmeans.png}
\caption{\emph{k-means} et mélange de Gaussiennes sur l'exemple de VanderPlas.}\label{fig:km}\end{figure}

\section{Conclusion}
EM fait croître la log-vraisemblance à chaque itération grâce à des étapes explicites. Il peut cependant s'arrêter
en un point selle, comme avec l'initialisation symétrique, et il devient très lent quand les composantes se
recouvrent. Sur ces données, toutes les initialisations aléatoires mènent au même maximum. La précision est limitée par
l'erreur statistique en $1/\sqrt n$, et le BIC permet de choisir $K$.

%\appendix
%\section{Dérivations (manuscrites)}
%\fig[0.9\textwidth]{derivations.jpg}

\begin{thebibliography}{9}
\bibitem{ee227c} M.~Hardt, EE~227C Lecture~19, UC Berkeley, 2018, \url{https://ee227c.github.io/code/lecture19.html\#em}.
\bibitem{ssbd} S.~Shalev-Shwartz, S.~Ben-David, \emph{Understanding Machine Learning}, Cambridge Univ.\ Press, 2014,
  ch.~24.
\bibitem{vdp} J.~VanderPlas, \emph{Python Data Science Handbook}, O'Reilly, 2016, §5.12.
\end{thebibliography}
\end{document}
